\section{Diseño de agentes autónomos}

\subsection{Definición y mecanismo de trabajo del agente}

A medida que los LLM se vuelven cada vez más maduros en capacidades clave —comprensión de entradas complejas, razonamiento y planificación, uso confiable de herramientas y recuperación de errores— los agentes comienzan a aparecer en entornos de producción.

El flujo de trabajo de un agente es el siguiente:
\begin{enumerate}
    \item Comienza con una orden del usuario humano o una discusión interactiva.
    \item Una vez definida la tarea, el agente planifica y opera de forma \textbf{independiente}.
    \item Durante la ejecución, obtiene «\textbf{Ground Truth}» del entorno mediante resultados de llamadas a herramientas, ejecución de código, etc., para evaluar el progreso.
    \item Puede pausarse en puntos de control (checkpoints) solicitando retroalimentación humana.
    \item Se detiene al completar la tarea; también suelen establecerse condiciones de parada como un número máximo de iteraciones.
\end{enumerate}

\begin{importantbox}{Esencia de la implementación del agente}
Un agente suena complejo, pero su implementación suele ser muy directa: \textbf{esencialmente consiste en que un LLM utilice herramientas basándose en retroalimentación del entorno dentro de un ciclo}. Lo clave no radica en una arquitectura compleja, sino en un conjunto de herramientas cuidadosamente diseñado y documentación clara de las mismas. La calidad del diseño de las herramientas determina el techo de rendimiento de un agente.
\end{importantbox}

\subsection{Cuándo utilizar un agente}

Los agentes son adecuados para los siguientes escenarios:
\begin{itemize}
    \item \textbf{Preguntas abiertas}: no se puede predecir el número de pasos necesarios.
    \item \textbf{No se puede codificar una ruta fija}: requiere que el modelo tome decisiones autónomas.
    \item \textbf{Entornos confiables}: existe cierto nivel de confianza en las decisiones del modelo.
    \item \textbf{Necesidades de escala}: requiere ejecutar tareas a gran escala dentro de un entorno de confianza.
\end{itemize}

\begin{warningbox}{Riesgos y costos de los agentes}
La autonomía de un agente implica: (1) \textbf{mayor costo}—el consumo de tokens por llamadas múltiples a herramientas supera ampliamente el de una llamada única; (2) \textbf{los errores pueden acumularse}—un error en una etapa puede amplificarse en etapas posteriores. Anthropic recomienda realizar pruebas exhaustivas en un \textbf{entorno de aislamiento} y establecer Guardrails apropiados (por ejemplo, número máximo de iteraciones, operaciones sensibles requieren confirmación humana, etc.).
\end{warningbox}

\subsection{Casos prácticos de agentes}

Dos implementaciones de agentes propias de Anthropic:

\begin{table}[H]
\centering
\begin{tabular}{>{\bfseries}p{3.5cm} p{9cm}}
\toprule
Caso & Descripción \\
\midrule
SWE-bench Coding Agent & Basado en la descripción de una pull request, edita automáticamente múltiples archivos para resolver un problema real de GitHub Issue. Puede resolver problemas reales en el benchmark SWE-bench Verified. \\
Computer Use Agent & Claude utiliza el equipo de cómputo (ratón, teclado, pantalla) para completar diversas tareas; es un caso típico de agente multimodal. \\
\bottomrule
\end{tabular}
\caption{Casos prácticos de implementación de agentes de Anthropic}
\end{table}

\subsection{Resumen de la sección}
Un agente representa el extremo superior del espectro de complejidad de los sistemas agenticos y es adecuado para tareas abiertas con rutas impredecibles. Su esencia de implementación es un ciclo de LLM más llamadas a herramientas; lo clave radica en el diseño de las herramientas, no en la complejidad arquitectónica. Al utilizarlos, hay que prestar atención al costo, al riesgo de acumulación de errores y a realizar pruebas exhaustivas en entornos de aislamiento.
